\section{Implementation and Evaluation}
\label{sec:eval}

The repository implements the full pipeline as a reproducible local artifact stack. Python ingestors validate and land provider exports, dbt models materialize staging, intermediate, and mart tables in DuckDB, and downstream Python modules consume the unified mart for allocation, intelligence, forecasting, and lifecycle persistence. The implementation is intentionally transparent rather than distributed: it favors inspectable data products and deterministic rule logic over opaque service orchestration.

\subsection{Benchmark and Reproducibility}
Evaluation is performed on a synthetic multi-cloud benchmark that exercises the financial and governance cases the paper claims to handle. The benchmark includes commitment-backed rows, shared-cost scenarios, tag gaps, unattributed spend, anomaly-style month-to-month variation, and lifecycle records. Reproducibility is supported directly in the repository through synthetic data generators, a dbt build path, a DuckDB-backed paper query pack, exported CSV and Markdown result sets, and LaTeX figure-generation scripts. Every headline number in the paper can therefore be regenerated from the checked-in pipeline rather than copied from one-off analysis.

\subsection{Current Benchmark Snapshot}
The current database state contains 18{,}992 normalized rows across four billing months. Aggregate list cost is \$67{,}065.52 and aggregate NEC is \$52{,}021.31, leaving \$15{,}044.21 of savings relative to list price and \$136.80 of explicitly modeled commitment waste. The benchmark also materializes \$29{,}874.07 of unattributed NEC, which serves as a governance stress case because it forces the system to distinguish visibility from accountability. On the decision side, the repository currently materializes 95 recommendations with \$1{,}709.83 of expected savings represented in lifecycle rows.

\subsection{Evidence Summary}
Several implemented invariants are already supported by current outputs. Shared-cost preservation holds exactly in the Azure benchmark path after fan-out, with zero delta between shared NEC and allocated shared NEC in each observed month. NEC summaries are internally consistent across cloud-level rollups and commitment-utilization measures. Recommendation outputs are auditable back to signal types, action mappings, and score components. Forecast quality is summarized by 68 backtest evaluations with \$7.86 mean absolute error, 2.98\% mean absolute percentage error, and a 95.59\% within-10-percent rate. The lifecycle layer is implemented end to end, but the benchmark should still be described as pre-realization: recommendations are persisted and trackable, yet realized-savings evidence remains sparse because current rows are concentrated in the \texttt{recommended} state.

\subsection{Scope and Limitations}
The evaluation is intentionally bounded. The benchmark is synthetic rather than enterprise production billing data. Waste detection relies on billing-derived proxies and commitment semantics rather than runtime telemetry. Recommendation logic is deterministic and auditable, but not autonomous remediation. These limitations do not undercut the systems contribution; they define it. The paper evaluates whether the implemented pipeline is coherent, reproducible, and internally validated as a post-billing FinOps decision engine, not whether it has already exhausted the broader space of production cloud cost optimization.
